% GENERATED by md2tex.py from ../draft_v5_en/body/11_discussion.md -- do not edit by hand.
% Change the Markdown and run `python md2tex.py`; edits here are overwritten.

\chapter{Discussion}
\label{ch:11}

\section{Purpose and method}
\label{sec:11.1}

This chapter evaluates the framework, not the system. This follows from \S\ref{sec:3.2.1}, where the work was presented as design science research: the object being judged is the adapted V-Model, and the lane-following system is the tool used to test it. The method was fixed in advance in \S\ref{sec:3.7}, with five criteria and their indicators, plus an evaluation of the three hypotheses from Chapter~\ref{ch:01}. Setting the criteria before knowing the results is what makes this evaluation more than a defence.

\section{Evaluation against the five criteria}
\label{sec:11.2}

\textbf{(1) Traceability integrity. Criterion: zero orphans. Met.} The checker found no orphans at any review gate or in the final matrix. What matters is not the final zero but the fact that the constraint worked as a hard gate during development: orphan identifiers were caught when they were added, not in a later audit. A zero obtained in a final audit would only show how careful the author was. A zero maintained continuously shows a property of the framework.

\textbf{(2) Requirement coverage by evidence. Criterion: 100\,\% with a verdict, even a negative one. Met.} All fourteen requirements have a verdict supported by quantitative evidence. The rule that an honest verdict is better than a missing one was applied in three ways. The global verdict of the camera arm was recorded as literally negative. Two requirements keep their literal failure next to the explanation instead of being rewritten. And one requirement is closed as not satisfied, with its measurement and a breakdown by rule combination. The criterion was met not in spite of the uncomfortable results but through them.

\textbf{(3) Hazard anticipation. Criterion: most of what was observed had been predicted, and what was not can be audited. Mostly met, and only partly tested by the physical step.} Almost every hazard that appeared was already in the register from the analysis phase, or was added when the camera track started, before the campaign measured it. Two cases were not predicted. The first is the conflict between the cage and its own estimator, where the envelope makes a competent policy worse because of a reading that is confident but wrong; its root cause was registered as a hazard, but not this way of appearing. The second is the degenerate optimum of stopping, which appears when the policy is given control of speed: the register predicted the family, reward exploitation, but the specific case had to be discovered. In both cases the framework handled the new case through the planned path of identifier, recorded decision and artefact.

A third unexpected case needs its own mention, because the framework did not predict it in any form: the dependence on a class B rule to stay in the lane (\S\ref{sec:8.5}). The hazard register considers that the cage may fail, arbitrate badly or read badly. It does not consider that the policy may adapt to the cage so much that it needs the cage in order to drive. This is a failure mode of the \emph{coupling} between the learned component and the envelope, and it suggests a hazard category that the original analysis did not have.

\textbf{The physical step tests this criterion more severely than simulation did.} On hardware the failures fall into two groups, which say different things. The first group is failures of families that were already registered but took a form the register did not describe. The register considers that the cage estimator may read badly and with high confidence, but it does not consider that its accuracy may depend on the place on the circuit (\S\ref{sec:9.3.6}). Because the family had been predicted, the case could be handled; because its form had not, it had to be measured before anything could be done. The second group is not about hazards at all, and the criterion covers it only by omission: rules that are correct according to their specification but have no defined \emph{behaviour in operation}, and thresholds that cannot be reached in the deployed configuration. A hazard register asks what can go wrong. It does not ask what happens next, and it does not ask whether the operating point ever reaches the envelope that was specified.

The honest reading is therefore that the criterion is met in simulation and only partly tested on hardware. No physical failure required a new hazard family, several required a new description, and the most uncomfortable type, which is the gap between specification and operation, lies outside what the tool can see. \S\ref{sec:9.5} measures the same point from the other side.

\textbf{(4) Adoption cost. Criterion: in proportion to the benefit. Met, provisionally.} The record holds more than seventy decisions, the checker, the templates and a dozen living documents. One person kept all of them up to date \emph{while} two training tracks, several evaluation campaigns and a bridge between simulators were being run. The exact split between framework work and technical work was not measured in hours, which was a limit stated in advance. The recorded indicator is that no review gate was delayed by framework artefacts, while some were delayed for technical reasons.

\textbf{(5) Usefulness of the matrix. Criterion: documented cases where it made the impact analysis faster. Met, with a counterexample.} There are three cases. The defect in the criterion that caused the literal negative verdict was found in minutes by following the row requirement $\rightarrow$ scenario $\rightarrow$ clause, because the chain existed. Re-scoring a requirement on its own metric was possible because the matrix showed which scenario tests which requirement with which metric. And moving the library to a new layout reused the existing chains, with changes that could be audited.

The counterexample is recorded with the same weight, because it shows a real limit. The matrix did not detect by itself that a scenario was being run without injecting its starting conditions. For some time that scenario produced zero co-activation and an indeterminate verdict that looked like an instrumentation gap. Traceability guarantees that the scenario exists and points to its requirement; it does not guarantee that the runner executes it as specified. That difference, checking the \emph{execution} rather than the \emph{reference}\sidenote{Checking the reference: every scenario ID points to a requirement that exists. Checking the execution: the scenario really injects its starting conditions and logs what its criterion needs.}, is the most concrete improvement this work identifies for its own framework.

\section{Evaluation of the three hypotheses}
\label{sec:11.3}

\textbf{H1 (construct): supported.} The five adaptations were enough to cover the failure modes found during the cycle, and none of them required a mechanism from outside the framework. The condition is that \enquote{enough} here means that nothing was found that did not fit. That is evidence that they are enough in practice, not that they are complete.

\textbf{H2 (operability): supported.} Each adaptation produced concrete artefacts, maintained by one person, and none became the bottleneck of the project. The cheapest adaptation was traceability, implemented as a checker of a few hundred lines, and it was also the one that gave the most value. This is the clearest cost-benefit result of the work.

\textbf{H3 (usefulness): supported, with an important condition.} The framework produced a well-founded verdict with its limits of validity, which is exactly what the hypothesis stated. The condition is that the verdict is uncomfortable: negative overall in its literal form, with one unmet requirement and one rule that was never tested. A reader could see this as a failure of the system, but that would mix two levels. The hypothesis is about whether the framework can produce a traceable judgement, not about whether that judgement is positive. A framework that could only produce positive judgements would not be useful.

\section{Lessons learned}
\label{sec:11.4}

\textbf{Traceability is cheap, instrumentation is expensive.} The checker cost little and gave a lot in return. What took time was making the scenarios measure what they claimed to measure: injecting starting conditions, logging the quantities the criteria use, and distinguishing indeterminate results from failures. The lesson for anyone reusing the framework is to plan the effort for instrumentation, not for documentation.

\textbf{An inherited criterion that is never audited contaminates the verdict.} A performance clause calibrated on an older controller and an older layout survived two migrations and blocked the global verdict of two campaigns. The audit showed that it measured something different from what its name suggested. Acceptance criteria are version-controlled artefacts and must be migrated as carefully as the code.

\textbf{Training metrics do not describe how a policy behaves.} Seeds with almost identical curves produced policies on opposite sides of the line between respecting the constraints and depending on the envelope. A framework that allowed selection by reward would have delivered, in this work, the worst of three candidate policies.

\textbf{A measurement at one spot does not describe an operating envelope.} This is the lesson the physical step pushed hardest, and the easiest one to repeat. The camera calibration and all the pre-deployment checks were done with the vehicle stopped at the same convenient place on the circuit. A capture over the whole layout then showed that this spot was the best one for the estimator (\S\ref{sec:9.3.6}). The content of those measurements still held; what was withdrawn was how widely they applied. The lesson is that a measurement protocol must cover the whole domain about which something will be claimed, and that an acceptance gate based on spread cannot detect a stable bias.

\textbf{The envelope can shape the component it contains.} This is the most unexpected lesson and the one with the widest consequences. When a policy is trained with the cage in the actuation chain, what is optimised is the pair, not the policy. Runtime containment is not a neutral filter placed on top of a finished component. It is part of the system in which the component learns to operate.

\section{Comparison with other practices}
\label{sec:11.5}

Compared with using no framework, the visible difference is that every claim has a place to go: each result has exactly one place where it is recorded, and none is left without evidence. Compared with model-based engineering, the framework gives up the power of a central model. In exchange it has a much lower entry barrier, since it needs only version-controlled text files and a checker, and it keeps the key property of verified traceability. Compared with pure safe RL, it keeps a guarantee that does not depend on the training distribution and that still holds if the policy is replaced. Compared with scenario-based validation without a life cycle, it adds the upward link: every scenario exists because a requirement asks for it and a hazard justifies it, not because it seemed interesting.

\section{Residual risks}
\label{sec:11.6}

Five risks remain open at the end, and they are stated without mitigation. Arbitration when rules fire together does not meet its requirement, and fixing it is an open design problem. The dependence on the rate limiter means that what transfers to hardware is a coupled pair and not a policy. The common cause between the policy's perception and the cage's perception is structural, because they share the image and a strong enough degradation blinds both at once; the existing mitigations fall back to a safe stop, but they do not remove the dependency. Physical evidence exists but is not scored: the chain drives on the real vehicle, but no run was done under the scenario protocol, so no verdict in this work is validated there. The fifth risk is the one that limits most what the framework can claim, and it is stated in \S\ref{sec:10.4}e: because the driving was always done in monitoring mode, the central artefact of this thesis, a cage that corrects actions at runtime, has been verified in simulation and only \emph{observed} on hardware. The framework produced the evidence that makes it possible to state this precisely, which is its job. It does not produce, and does not claim to produce, the evidence needed to state the opposite.

\section{Summary}
\label{sec:11.7}

The framework meets the five criteria it set for itself, including the one that required verdicts for all requirements even when they were unfavourable. Its most useful limit is not one of those stated in advance but the one it revealed when it was applied: traceability checks references, not executions, and that gap allowed one scenario to run for months without doing what its specification said. Chapter~\ref{ch:12} collects the conclusions and turns these findings into concrete lines of future work.
